\documentclass[11pt]{article}
\usepackage{amsmath,amssymb,amsthm}
\newtheorem{theorem}{Theorem}
\title{Three compact reverse heads shatter eight columns}
\author{}
\date{}
\begin{document}
\maketitle

\paragraph{Status and scope.}
This is a complete written proof, supplemented by an exact rational
checker. It has not been formalized in Lean. The improvement is a fixed
local constant, not either accepted asymptotic target. No growing
query-to-head ratio follows from this argument. The construction initially
has no axis-zero requirement; the existing checked axis repair can then
be applied with its stated constant overhead.

\begin{theorem}
For every integer $r\geq1$ and every sign table
$\sigma:[r]\times[8]\to\{-1,1\}$, three affine ratios realize the
strict grid signs $\sigma$. Every denominator has strictly negative
slopes and is positive on the entire Boolean $(r+8)$-cube.
\end{theorem}

\begin{proof}
We first construct eight fixed normalized queries for the three-head
family
\[
 f_j(a,p)=\sum_{h=1}^3\frac{a_h+b_{hj}}{D_{hj}-p_h}.
\]
Put $t=101/100$, $a^0=(1,1,-2)$, and
$w^A=(1,1,1)$, $w^B=(101/100,99/100,1)$.
The first four queries have zero $b$ and respective denominator vectors
$1/w^A,t/w^A,1/w^B,t/w^B$, where division is coordinatewise.
For the remaining four queries use the following pairs of vectors $w,v$:
\[
\begin{array}{c|ccc|ccc}
 &w_1&w_2&w_3&v_1&v_2&v_3\\ \hline
1&101/100&24/25&1&10&-8&-475/303\\
2&51/50&26/25&49/50&-5&-10&18865/1326\\
3&1&51/50&101/100&7&7&-71407/5100\\
4&1&99/100&1&0&-2&200/99
\end{array}
\]
Set $D_h=1/w_h$ and $b_h=v_h/w_h^2$ for each such query.
By construction $\sum_hv_h/w_h=0$, so all eight queries vanish
at $a=p=0$.

Restrict the first-order numerator direction to $a=\delta\lambda a^0$,
and write $p=\delta P$. The four extra query derivatives are
$L(\lambda,P_1,P_2,P_3)^{\mathsf T}$, where
\[
L=\begin{pmatrix}
-3/100&10&-8&-475/303\\
1/10&-5&-10&18865/1326\\
0&7&7&-71407/5100\\
-1/100&0&-2&200/99
\end{pmatrix},\qquad
\det L=\frac{835605281}{736593000}>0.
\]
Given $Y\in\mathbb R^4$, set $(\lambda,P)=L^{-1}Y$ and define
\[
 Q_A=\lambda(P_1+P_2-2P_3),\qquad
 Q_B=\lambda((101/100)^2P_1+(99/100)^2P_2-2P_3).
\]
For completeness, the following exact rows give the three linear forms
$\lambda=uY$, $P_1+P_2-2P_3=v_A Y$, and
$(101/100)^2P_1+(99/100)^2P_2-2P_3=v_B Y$:
\[
\begin{array}{c|rrrr}
u&3643867900/358116549&371958600/119372183&-30859300000/2506815843&-11861500200/119372183\\
v_A&72215/238744366&-86658/596860915&118924450/835605281&-563277/238744366\\
v_B&2906209249/1432466196000&1868678071/2984304575000&71834075071/501363168600&418378323/2387443660000
\end{array}
\]
These identities can also be checked by multiplying the displayed rows
by $L$. Every entry is nonzero. For each of these three rows $v$ and
every coordinate $j$, direct rational comparison gives
\[
 \sum_{k\ne j}|v_k|<1000|v_j|.
\]
The diagonal sign pairs $(\operatorname{sgn}(u_jv_{A,j}),
\operatorname{sgn}(u_jv_{B,j}))$, in order, are
$(+,+),(-,+),(-,-),(+,-)$.

Now prescribe all eight query signs. Choose the coordinate $j$ whose
diagonal sign pair is the desired sign at the first $A$ query and the
first $B$ query. Set $Y_j$ equal to $1000$ times its desired extra-query
sign, and each other $Y_k$ equal to its desired extra-query sign.
The strict dominance inequalities show that $Q_A,Q_B$ have the two
prescribed signs, independently of the four extra-query signs.

For each pair $K=A,B$, let its two desired signs be $s_K,s'_K$ and set
\[
 \ell_K=\begin{cases}
 0&s_K=s'_K,\\
 -Q_K(t+1)/(2t)&s_K\ne s'_K.
 \end{cases}
\]
Choose
\[
 a^2=(50(\ell_B-\ell_A),-50(\ell_B-\ell_A),\ell_A).
\]
Then $w^A\cdot a^2=\ell_A$ and $w^B\cdot a^2=\ell_B$.
Take the actual normalized parameters
\[
 a=\delta\lambda a^0+\delta^2a^2,\qquad p=\delta P.
\]
The first-order terms in each pair vanish, since
$w^A\cdot a^0=w^B\cdot a^0=0$. The two scores of pair $K$ are
\[
 \delta^2(\ell_K+Q_K)+O(\delta^3),\qquad
 \frac{\delta^2}{t^2}(t\ell_K+Q_K)+O(\delta^3).
\]
If the signs agree, the leading terms both have sign $Q_K$.
If they disagree, the parenthesized leading terms are respectively
$Q_K(t-1)/(2t)$ and $-Q_K(t-1)/2$. All are nonzero and have the
required signs. Each extra query is $\delta Y_k+O(\delta^2)$,
and therefore has its required sign as well. All normalized denominators
are positive for sufficiently small positive $\delta$.

There are only $256$ desired sign vectors. The preceding construction
provides fixed rational $\lambda,P,a^2$ for every one, so one common
positive rational $\delta$ can be chosen to work for all of them.
This assertion is constructive: clear the three denominators of each
query. Its numerator is a polynomial of degree at most four in $\delta$,
whose first nonzero coefficient has the proved desired sign. If this
coefficient is $c_k$, then
\[
0<\delta\leq\min\left(\frac12,
 \frac{|c_k|}{2\sum_{l>k}|c_l|}\right)
\]
preserves the sign, omitting the second bound when its denominator is
zero. Take the minimum over all queries and all $256$ patterns, together
with the positive denominator bounds. This is a finite rational
algorithm, with no numerical feasibility hypothesis.

It remains to put arbitrarily many rows on one admissible cube. All
normalized query denominators satisfy
$25/26\leq D_{hj}\leq25/24$. Put
$d_0=19/20$ and $\rho=1/(1000r)$. Choose the preceding common
$\delta$ smaller, if necessary, so that
$d_0\delta|P_h|<\rho/2$ for every pattern and head.
In row $i$ use the parameters for its desired eight-sign vector and set
\[
 \widehat p_{hi}=\rho+d_0\delta P_{hi},\qquad
 \widehat q_{hj}=1-\rho-d_0D_{hj}.
\]
Define the three affine denominators and numerators by
\[
 B_h(z,w)=1-\sum_i\widehat p_{hi}z_i-\sum_j\widehat q_{hj}w_j,
 \qquad
 A_h(z,w)=d_0\left(\sum_i a_{hi}z_i+\sum_jb_{hj}w_j\right).
\]
At $(e_i,e_j)$ the ratio equals its normalized counterpart, since
$B_h(e_i,e_j)=d_0(D_{hj}-\delta P_{hi})$.
Moreover
\[
 \rho/2<\widehat p_{hi}<3\rho/2,
 \qquad 0<1/96-\rho\leq\widehat q_{hj}<9/104.
\]
The total slope magnitude is consequently less than
$3r\rho/2+8(9/104)=3/2000+9/13<1$.
All slopes are strictly negative, and the minimum denominator on the
entire Boolean cube is positive. This completes the proof.
\end{proof}

\paragraph{Checked transfer and remaining limitation.}
The Lean theorem
\texttt{hasReverseAxisCompiler\_of\_strictReverseGridCompiler}
in \texttt{Results/ReverseGridAxisRepair.lean} converts an $H$-head
strict reverse grid compiler to an axis-zero compiler with $3H+3$ heads.
Thus this written construction gives twelve axis-zero heads for every
$r\times8$ table. Its use with this written input theorem has not been
formalized. The construction exhibits two independent quadratic
disagreement choices in addition to four linear choices. It supplies no
induction producing an unbounded query-to-head ratio, so neither the
global logarithmic saving nor a matching global lower bound is proved.
\end{document}
